%! TEX root = master.tex
\chapter{2}{Conics and Their Reflective Properties}

Parabolas and ellipses in standard position. The tools are implicit
differentiation for tangent slopes and the angle formula for comparing two
lines. The chapter ends with the reflective property of the ellipse.

\section{The parabola}\label{sec:parabola}

%! TEX root = ../master.tex
\begin{theorem}[Parabola in standard position]\label{thm:parabola-standard}
  The parabola with focus $(h + p, k)$ and directrix $x = h - p$, where
  $p \neq 0$, is
  \[
    (y - k)^2 = 4p(x - h).
  \]
  Its vertex is $(h, k)$ and its axis is the line $y = k$. Exchanging the roles
  of $x$ and $y$ gives the parabola $(x - h)^2 = 4p(y - k)$ with a vertical axis.
\end{theorem}


\begin{exercise}\label{ex:parabola-tangents}
  Tangent lines $T_1$ and $T_2$ to the parabola $y = x^2$ at the points $P_1$
  and $P_2$ meet at $P$. A third tangent $T$, touching the parabola at a
  point between $P_1$ and $P_2$, meets $T_1$ at $Q_1$ and $T_2$ at $Q_2$.
  Show that
  \[
    \frac{PQ_1}{PP_1} + \frac{PQ_2}{PP_2} = 1.
  \]
\end{exercise}

\begin{figure}[htbp]
  \centering
  \includegraphics[width=0.62\linewidth]{parabola-tangents}
  \caption{Three tangents to $y = x^2$.}
  \label{fig:parabola-tangents}
\end{figure}

\begin{solution}
  Let $P_1 = (a, a^2)$ and $P_2 = (b, b^2)$, and let $T$ touch at
  $(c, c^2)$ with $c$ strictly between $a$ and $b$. In
  \cref{fig:parabola-tangents} we have $b < c < a$, but the argument below
  does not depend on the order.

  \emph{Tangents.} Since $\frac{\dd y}{\dd x} = 2x$, the tangent at
  $(s, s^2)$ is $y - s^2 = 2s(x - s)$, that is, $y = 2sx - s^2$. So
  \[
    T_1 : y = 2ax - a^2, \qquad
    T_2 : y = 2bx - b^2, \qquad
    T : y = 2cx - c^2.
  \]
  \emph{Intersections.} Tangents at $s \neq u$ meet where
  $2(s - u)\,x = s^2 - u^2$, that is, at $x = \frac{s + u}{2}$, and then
  $y = 2s \cdot \frac{s + u}{2} - s^2 = su$. Hence
  \[
    P\Bigl(\frac{a + b}{2}, ab\Bigr), \qquad
    Q_1\Bigl(\frac{a + c}{2}, ac\Bigr), \qquad
    Q_2\Bigl(\frac{b + c}{2}, bc\Bigr).
  \]
  \emph{Ratios.} The points $P$, $Q_1$, $P_1$ lie on $T_1$, which is not
  vertical. By \cref{thm:section-formula}(b), the ratio of lengths equals the
  ratio of changes in $x$:
  \[
    \frac{PQ_1}{PP_1}
    = \frac{\abs{\frac{a + b}{2} - \frac{a + c}{2}}}{\abs{\frac{a + b}{2} - a}}
    = \frac{\abs{b - c}}{\abs{b - a}}.
  \]
  In the same way, on $T_2$,
  \[
    \frac{PQ_2}{PP_2}
    = \frac{\abs{\frac{b + c}{2} - \frac{a + b}{2}}}{\abs{b - \frac{a + b}{2}}}
    = \frac{\abs{c - a}}{\abs{b - a}}.
  \]
  Because $c$ lies between $a$ and $b$, $\abs{b - c} + \abs{c - a} = \abs{b - a}$,
  and so
  \[
    \frac{PQ_1}{PP_1} + \frac{PQ_2}{PP_2}
    = \frac{\abs{b - c} + \abs{c - a}}{\abs{b - a}} = 1. \qedhere
  \]
\end{solution}

\begin{remark}
  The intersection rule is worth remembering: the tangents at $s$ and $u$
  meet at $\bigl(\frac{s + u}{2}, su\bigr)$, above the midpoint of $s$ and
  $u$.
\end{remark}

\begin{exercise}\label{ex:parabola-normal}
  Let $P(x_1, y_1)$ be a point other than the vertex on the parabola
  $y^2 = 2px$. Find the normal line at $P$, and show that it meets the
  $x$-axis at $Q(x_1 + p, 0)$.
\end{exercise}

\begin{figure}[htbp]
  \centering
  \includegraphics[width=0.45\linewidth]{parabola-normal}
  \caption{The normal at $P$. The label on $Q$ should read
    $(x_1 + p, 0)$.}
  \label{fig:parabola-normal}
\end{figure}

\begin{solution}
  Differentiating $y^2 = 2px$ implicitly gives $2y\,\frac{\dd y}{\dd x} = 2p$,
  so the tangent at $P$ has slope $p / y_1$. Here $y_1 \neq 0$ because $P$ is
  not the vertex. By \cref{thm:perpendicular-slopes}, the normal has slope
  $-y_1 / p$, and its equation is
  \[
    y - y_1 = -\frac{y_1}{p}(x - x_1).
  \]
  At the $x$-axis $y = 0$, so $-y_1 = -\frac{y_1}{p}(x - x_1)$. Dividing by
  $-y_1 \neq 0$ gives $1 = \frac{x - x_1}{p}$, that is, $x = x_1 + p$. So the
  normal meets the axis at $Q(x_1 + p, 0)$.
\end{solution}

\begin{remark}
  This is the reflective property of the parabola in disguise. Writing
  $y^2 = 2px$ as $y^2 = 4 \cdot \frac p2 \cdot x$, \cref{thm:parabola-standard}
  puts the focus at $F(\frac p2, 0)$ and the directrix at $x = -\frac p2$.
  By the focus--directrix definition,
  $FP = x_1 + \frac p2 = FQ$, so the triangle $FPQ$ is isosceles. The normal
  $PQ$ therefore makes equal angles with $PF$ and with the horizontal line
  through $P$. In other words, rays parallel to the axis reflect through the
  focus.
\end{remark}

\section{The ellipse}\label{sec:ellipse}

%! TEX root = ../master.tex
\begin{theorem}[Ellipse in standard position]\label{thm:ellipse-standard}
  Let $0 \leq c < a$ and $b^2 = a^2 - c^2$. The set of points whose distances
  to the foci $F_1(-c, 0)$ and $F_2(c, 0)$ add up to $2a$ is the ellipse
  \[
    \frac{x^2}{a^2} + \frac{y^2}{b^2} = 1.
  \]
  It meets the $x$-axis at $(\pm a, 0)$ and the $y$-axis at $(0, \pm b)$.
\end{theorem}


\begin{exercise}\label{ex:overpass}
  A highway overpass has the shape of half an ellipse, with a base 200 feet
  wide. The highway beneath it is 120 feet wide, and the minimum clearance
  over the highway is 40 feet. Find the height of the overpass at its
  centre.
\end{exercise}

\begin{figure}[htbp]
  \centering
  \includegraphics[width=0.55\linewidth]{elliptical-overpass}
  \caption{The overpass as the upper half of
    $\frac{x^2}{a^2} + \frac{y^2}{b^2} = 1$.}
  \label{fig:overpass}
\end{figure}

\begin{solution}
  Put the centre of the base at the origin, so the arch is the upper half of
  $\frac{x^2}{a^2} + \frac{y^2}{b^2} = 1$ and the height at the centre is $b$.

  \emph{Base.} The base is the major axis, so $2a = 200$ and $a = 100$.

  \emph{Clearance.} The highway runs from $x = -60$ to $x = 60$. On the
  upper half, $y$ decreases as $\abs{x}$ grows, so the clearance over the road
  is smallest at its edges. Therefore $(60, 40)$ lies on the ellipse:
  \[
    \frac{60^2}{100^2} + \frac{40^2}{b^2} = 1
    \implies \frac{1600}{b^2} = 1 - 0.36 = 0.64
    \implies b^2 = 2500.
  \]
  Since $b > 0$, $b = 50$. The overpass is $50$ feet high at its centre.
\end{solution}

\section{Reflection}\label{sec:reflection}

%! TEX root = ../master.tex
\begin{theorem}[Angle between two lines]\label{thm:angle-between-lines}
  Let $l_1$ and $l_2$ be non-perpendicular lines with slopes $m_1$ and $m_2$,
  and let $\alpha$ be the angle from $l_1$ to $l_2$ (measured
  anticlockwise). Then
  \[
    \tan \alpha = \frac{m_2 - m_1}{1 + m_1 m_2}.
  \]
\end{theorem}


\begin{figure}[htbp]
  \centering
  \includegraphics[width=0.7\linewidth]{ellipse-reflection}
  \caption{A ray from $F_1$ meets the ellipse at $P$ and reflects to $F_2$.
    The tangent at $P$ is dashed.}
  \label{fig:ellipse-reflection}
\end{figure}

%! TEX root = ../master.tex
\begin{theorem}[Reflective property of the ellipse]\label{thm:ellipse-reflection}
  A ray leaving one focus of an ellipse is reflected by the ellipse through
  the other focus.
\end{theorem}

\begin{proof}
  Take the ellipse of \cref{thm:ellipse-standard} with $c > 0$ (if $c = 0$
  it is a circle, and every ray from the centre returns along itself). Let
  $P(x_1, y_1)$ be on the ellipse.

  If $y_1 = 0$, then $P = (\pm a, 0)$, the tangent there is vertical, and the
  ray runs along the $x$-axis, so it returns along the axis through both
  foci. Suppose now $y_1 \neq 0$ and, for the moment, $x_1 \neq \pm c$, so
  that neither focal line is vertical.

  \emph{Slopes.} Differentiating $x^2/a^2 + y^2/b^2 = 1$ implicitly gives
  $2x/a^2 + (2y/b^2)\,y' = 0$. So the tangent at $P$ has slope
  \[
    m = -\frac{b^2 x_1}{a^2 y_1},
  \]
  and the lines $F_1P$ and $F_2P$ have slopes
  \[
    m_1 = \frac{y_1}{x_1 + c}, \qquad m_2 = \frac{y_1}{x_1 - c}.
  \]
  Below we use $b^2 x_1^2 + a^2 y_1^2 = a^2 b^2$ (because $P$ is on the
  ellipse) and $a^2 - b^2 = c^2$.

  \emph{Incoming angle.} Let $\alpha$ be the angle from $F_1P$ to the tangent.
  By \cref{thm:angle-between-lines},
  $\tan\alpha = (m - m_1)/(1 + m m_1)$. Multiply the numerator and the
  denominator by $a^2 y_1 (x_1 + c)$:
  \[
    \begin{aligned}
      \text{numerator}
      &= -b^2 x_1 (x_1 + c) - a^2 y_1^2
       = -(b^2 x_1^2 + a^2 y_1^2) - b^2 c x_1
       = -b^2 (a^2 + c x_1), \\
      \text{denominator}
      &= a^2 y_1 (x_1 + c) - b^2 x_1 y_1
       = y_1 \bigl( (a^2 - b^2) x_1 + a^2 c \bigr)
       = c y_1 (a^2 + c x_1).
    \end{aligned}
  \]
  Since $\abs{x_1} \leq a$ and $c < a$, we have $a^2 + c x_1 > 0$. So the
  denominator is not zero (the lines are not perpendicular), and
  \[
    \tan\alpha = -\frac{b^2}{c y_1}.
  \]

  \emph{Outgoing angle.} Let $\beta$ be the angle from the tangent to $F_2P$,
  so $\tan\beta = (m_2 - m)/(1 + m m_2)$. Multiply the numerator and the
  denominator by $a^2 y_1 (x_1 - c)$:
  \[
    \begin{aligned}
      \text{numerator}
      &= a^2 y_1^2 + b^2 x_1 (x_1 - c)
       = b^2 (a^2 - c x_1), \\
      \text{denominator}
      &= a^2 y_1 (x_1 - c) - b^2 x_1 y_1
       = -c y_1 (a^2 - c x_1),
    \end{aligned}
  \]
  with $a^2 - c x_1 > 0$. Hence
  \[
    \tan\beta = -\frac{b^2}{c y_1} = \tan\alpha.
  \]

  Both angles lie in $[0, \pi)$, where $\tan$ is one-to-one, so
  $\alpha = \beta$. So $F_1P$ and $F_2P$ make equal angles with the tangent,
  one on each side: reflecting the line $F_1P$ in the tangent gives the line
  $F_2P$. This is the law of reflection (the angle of incidence equals the
  angle of reflection), so the reflected ray passes through $F_2$.

  When $x_1 = \pm c$, one focal line is vertical and the slope formula does
  not apply. The result still holds there by continuity: both angles depend
  continuously on $P$, and they agree at every nearby point of the ellipse.
\end{proof}


\begin{remark}
  The simplifications in the proof rely on two facts: $P$ lies on the
  ellipse, so $b^2x_1^2 + a^2y_1^2 = a^2b^2$, and $a^2 - b^2 = c^2$. Without
  substituting both, the expressions for $\tan\alpha$ and $\tan\beta$ do not
  visibly agree.
\end{remark}
